\section*{第 6 章}
\addcontentsline{toc}{section}{第 6 章}

\exercise{习题 6.1}
\textbf{题目：}
设二进制序列中各符号相互统计独立，两个二进制符号等概率出现，信息速率
\[
R_b=1\,\mathrm{Mbit/s}.
\]
画出下列随机信号的平均功率谱密度图，并标上频率值：
\begin{enumerate}
\item 单极性矩形不归零码序列；
\item 该单极性矩形不归零码序列乘以载波 $\cos(2\pi f_ct)$ 后得到的 OOK 信号，其中 $f_c=100\,\mathrm{MHz}$。
\end{enumerate}

\par\textbf{解：}
比特间隔为
\[
T_b=\frac{1}{R_b}=1\,\mu s.
\]

\textbf{(1) 单极性 NRZ}\par
单极性 NRZ 可看成 $a_n\in\{0,1\}$、等概率独立的 PAM 信号，因此
\[
\mathbb{E}[a_n]=\frac12,\qquad \sigma_a^2=\frac14.
\]
矩形脉冲的频谱为
\[
G(f)=T_b\sinc(fT_b).
\]
双边平均功率谱密度为
\[
P_{NRZ}(f)=\frac{T_b}{4}\sinc^2(fT_b)+\frac14\delta(f).
\]
连续谱主瓣第一零点在
\[
f=\pm \frac{1}{T_b}=\pm 1\,\mathrm{MHz}.
\]

\textbf{(2) OOK 信号}\par
OOK 信号可写成
\[
s(t)=A\,m(t)\cos(2\pi f_ct).
\]
其双边平均功率谱密度为
\[
P_{OOK}(f)=\frac{A^2}{4}\Bigl[P_{NRZ}(f-f_c)+P_{NRZ}(f+f_c)\Bigr].
\]
代入上式可得
\[
P_{OOK}(f)=\frac{A^2T_b}{16}\Bigl[\sinc^2((f-f_c)T_b)+\sinc^2((f+f_c)T_b)\Bigr]
+\frac{A^2}{16}\bigl[\delta(f-f_c)+\delta(f+f_c)\bigr].
\]
因此频谱以 $\pm100\,\mathrm{MHz}$ 为中心，主瓣第一零点位于
\[
f=\pm(100\pm 1)\,\mathrm{MHz}.
\]

\medskip
\exercise{习题 6.2}
\textbf{题目：}
某 OOK 系统在区间 $[0,T_b]$ 内等概率发送
\[
s_1(t)=\sin\frac{4\pi t}{T_b},
\qquad
s_2(t)=0.
\]
OOK 信号通过 AWGN 信道传输，信道噪声双边功率谱密度为 $N_0/2$。接收端利用带通型匹配滤波器进行解调。

\begin{enumerate}
\item 画出匹配滤波器的冲激响应；
\item 画出匹配滤波器输出的有用信号波形；
\item 确定最佳判决门限，并求系统平均误比特率。
\end{enumerate}

\par\textbf{解：}
\textbf{(1) 匹配滤波器}\par
匹配滤波器应与 $s_1(t)$ 匹配，因此
\[
h(t)=s_1(T_b-t)
=\sin\!\left(4\pi-\frac{4\pi t}{T_b}\right)
=-\sin\frac{4\pi t}{T_b},
\qquad 0\le t\le T_b,
\]
其他时刻为 $0$。

\textbf{(2) 输出有用信号}\par
匹配滤波器输出的有用信号是
\[
y_s(t)=s_1(t)*h(t),
\]
也就是 $s_1(t)$ 的自相关波形。其峰值出现在最佳采样时刻 $t=T_b$。由于
\[
E_1=\int_0^{T_b}\sin^2\frac{4\pi t}{T_b}\,dt=\frac{T_b}{2},
\]
故峰值为
\[
y_s(T_b)=E_1=\frac{T_b}{2}.
\]
因此输出有用信号是一个在 $t=T_b$ 取最大值的相关峰。

\textbf{(3) 门限与误比特率}\par
由于“1”“0”等概率发送，最佳判决门限取两采样均值的中点：
\[
V_T=\frac{E_1+0}{2}=\frac{T_b}{4}.
\]
匹配滤波器采样噪声方差为
\[
\sigma^2=\frac{N_0}{2}E_1=\frac{N_0T_b}{4}.
\]
平均比特能量为
\[
E_b=\frac{E_1}{2}=\frac{T_b}{4}.
\]
因此平均误比特率为
\[
P_b=\frac12\operatorname{erfc}\!\left(\frac{V_T}{\sqrt{2\sigma^2}}\right)
=\frac12\operatorname{erfc}\!\left(\sqrt{\frac{E_b}{2N_0}}\right).
\]

\medskip
\exercise{习题 6.3}
\textbf{题目：}
已知二进制 2FSK 通信系统的两个信号波形为
\[
s_1(t)=\sin\frac{2\pi}{T_b}t,\qquad
s_2(t)=\sin\frac{4\pi}{T_b}t,
\qquad 0\le t\le T_b,
\]
其中 $T_b=1\,s$。信号在信道传输中受到均值为 $0$、双边功率谱密度为 $N_0/2$ 的 AWGN 干扰，且 $s_1(t)$ 与 $s_2(t)$ 等概率出现。求：
\begin{enumerate}
\item 两信号波形图；
\item 相关系数 $\rho$ 及平均比特能量 $E_b$；
\item 最佳接收框图；
\item 匹配滤波器冲激响应；
\item 最佳判决门限；
\item 平均误比特率。
\end{enumerate}

\par\textbf{解：}
\textbf{(1)} $s_1(t)$ 在一个比特间隔内完成 1 个周期，$s_2(t)$ 完成 2 个周期。

\textbf{(2)} 两信号的互相关为
\[
\int_0^{T_b}s_1(t)s_2(t)\,dt=0,
\]
故
\[
\rho=0.
\]
两信号能量相同：
\[
E_1=E_2=\int_0^{T_b}\sin^2\frac{2\pi kt}{T_b}\,dt=\frac{T_b}{2}=\frac12.
\]
故平均比特能量为
\[
E_b=\frac12.
\]

\textbf{(3)} 最佳接收机为两路并行带通匹配滤波器或相关器，分别与 $s_1(t)$、$s_2(t)$ 匹配，比较两路采样值大小后作判决。

\textbf{(4)} 两个匹配滤波器的冲激响应为
\[
h_1(t)=s_1(T_b-t)= -\sin\frac{2\pi t}{T_b},\qquad 0\le t\le T_b,
\]
\[
h_2(t)=s_2(T_b-t)= -\sin\frac{4\pi t}{T_b},\qquad 0\le t\le T_b.
\]

\textbf{(5)} 由于两信号等能量且等概率，最佳判决门限为
\[
V_T=0.
\]

\textbf{(6)} 若发送 $s_1(t)$，则判决统计量可写成
\[
y=E_b+Z,
\]
其中 $Z$ 为零均值高斯变量，方差为
\[
\sigma_Z^2=N_0E_b=\frac{N_0}{2}.
\]
因此平均误比特率为
\[
P_b=\frac12\operatorname{erfc}\!\left(\sqrt{\frac{E_b}{2N_0}}\right).
\]
由于 $E_b=\frac12$，也可写成
\[
P_b=\frac12\operatorname{erfc}\!\left(\frac{1}{2\sqrt{N_0}}\right).
\]

\medskip
\exercise{习题 6.4}
\textbf{题目：}
设二进制信息速率为 $1\,\mathrm{Mbit/s}$，二进制序列中两个二进制符号等概率出现且相互独立。画出下列随机信号的双边平均功率谱密度图，并标上频率值：
\begin{enumerate}
\item 双极性矩形不归零码序列；
\item 双极性矩形不归零码序列通过乘法器后的 BPSK 信号，设载频 $f_c=100\,\mathrm{MHz}$。
\end{enumerate}

\par\textbf{解：}
仍有
\[
T_b=1\,\mu s.
\]

\textbf{(1) 双极性 NRZ}\par
令 $a_n=\pm1$ 等概独立，则
\[
\mathbb{E}[a_n]=0,\qquad \mathbb{E}[a_n^2]=1.
\]
故双边平均功率谱密度为
\[
P_{bNRZ}(f)=T_b\sinc^2(fT_b).
\]
第一零点在
\[
f=\pm1\,\mathrm{MHz}.
\]

\textbf{(2) BPSK}\par
BPSK 信号可写成
\[
s(t)=A\,m(t)\cos(2\pi f_ct).
\]
其双边平均功率谱密度为
\[
P_{BPSK}(f)=\frac{A^2}{4}\Bigl[P_{bNRZ}(f-f_c)+P_{bNRZ}(f+f_c)\Bigr].
\]
即
\[
P_{BPSK}(f)=\frac{A^2T_b}{4}
\Bigl[
\sinc^2((f-f_c)T_b)+\sinc^2((f+f_c)T_b)
\Bigr].
\]
因此双边频谱关于 $0$ 对称，在 $\pm100\,\mathrm{MHz}$ 处为中心，第一零点位于
\[
f=\pm(100\pm1)\,\mathrm{MHz}.
\]
由于基带均值为零，BPSK 频谱中没有离散载波谱线。

\medskip
\exercise{习题 6.5}
\textbf{题目：}
BPSK 信号
\[
s_1(t)=A\cos\omega_ct,\qquad
s_2(t)=-A\cos\omega_ct,\qquad 0\le t\le T_b
\]
分别表示发送“1”和“0”。信道中叠加均值为 $0$、双边功率谱密度为 $N_0/2$ 的 AWGN，且两信号等概率出现。

求：
\begin{enumerate}
\item 相关系数 $\rho$ 及平均比特能量 $E_b$；
\item 利用相关型解调实现 BPSK 最佳接收的框图；
\item 发 $s_2(t)$ 时，最佳采样值 $y$ 的条件均值、条件方差及条件概率密度；
\item 平均误比特率。
\end{enumerate}

\par\textbf{解：}
\textbf{(1)} 由于
\[
s_2(t)=-s_1(t),
\]
故
\[
\rho=-1.
\]
平均比特能量为
\[
E_b=\int_0^{T_b}A^2\cos^2\omega_ct\,dt=\frac{A^2T_b}{2}.
\]

\textbf{(2)} 最佳相关接收机可表述为：接收信号与本地载波 $s_1(t)$ 做相关，采样后与零门限比较完成判决。

\textbf{(3)} 发 $s_2(t)$ 时，
\[
y=\int_0^{T_b}[s_2(t)+n(t)]s_1(t)\,dt=-E_b+Z,
\]
其中
\[
Z\sim\mathcal{N}\!\left(0,\frac{N_0E_b}{2}\right).
\]
故
\[
\mathbb{E}[y|s_2]=-E_b,
\qquad
D(y|s_2)=\frac{N_0E_b}{2}.
\]
条件概率密度为
\[
p_2(y)=\frac{1}{\sqrt{\pi N_0E_b}}
\exp\!\left[-\frac{(y+E_b)^2}{N_0E_b}\right].
\]

\textbf{(4)} 最佳门限为 $0$，故平均误比特率为
\[
P_b=\frac12\operatorname{erfc}\!\left(\sqrt{\frac{E_b}{N_0}}\right).
\]

\medskip
\exercise{习题 6.6}
\textbf{题目：}
在 AWGN 干扰下对 BPSK 信号进行相干解调，恢复载波与发送载波的相位差为固定常数 $\theta$，且 $|\theta|<\pi/2$。设 BPSK 采用矩形脉冲成形，已调信号幅度为 $A$。
\begin{enumerate}
\item 写出宽带白高斯噪声通过窄带带通滤波器后的噪声 $n(t)$ 的数学表达式，画出其功率谱密度图，并求均值与方差；
\item 证明该系统的平均误比特率为
\[
P_b=\frac12\operatorname{erfc}\!\left(\sqrt{\frac{E_b}{N_0}}\cos\theta\right).
\]
\end{enumerate}

\par\textbf{解：}
\textbf{(1)} 窄带带通噪声可表示为
\[
n(t)=n_c(t)\cos(2\pi f_ct+\theta)-n_s(t)\sin(2\pi f_ct+\theta),
\]
其中 $n_c(t)$、$n_s(t)$ 是相互独立的低通高斯过程。其均值为
\[
\mathbb{E}[n(t)]=0.
\]
功率谱密度只分布在 $\pm f_c$ 附近，带宽等于带通滤波器带宽 $B$，总方差为
\[
\sigma_n^2=\mathbb{E}[n^2(t)]=N_0B.
\]

\textbf{(2)} 相干解调后的采样统计量为
\[
y=\pm A\cos\theta+Z,
\]
其中
\[
Z\sim\mathcal{N}(0,\sigma^2).
\]
由矩形脉冲 BPSK 的等效关系可得
\[
\frac{A^2}{2}=E_b.
\]
仍以零为门限判决，则平均误比特率为
\[
P_b=P(Z>A\cos\theta)
=\frac12\operatorname{erfc}\!\left(\frac{A\cos\theta}{\sqrt{2\sigma^2}}\right)
=\frac12\operatorname{erfc}\!\left(\sqrt{\frac{E_b}{N_0}}\cos\theta\right).
\]

\medskip
\exercise{习题 6.7}
\textbf{题目：}
BPSK 信号
\[
s_1(t)=A\cos\omega_ct,\qquad
s_2(t)=-A\cos\omega_ct,\qquad 0\le t\le T_b
\]
中，二进制序列为不相关序列，$s_1(t)$ 与 $s_2(t)$ 等概率出现。
\begin{enumerate}
\item 画出该 BPSK 信号的双边平均功率谱密度图；
\item 画出 BPSK 接收机提取载波的两种方框图，并简要说明工作原理。
\end{enumerate}

\par\textbf{解：}
\textbf{(1)} 由双极性 NRZ 调制得到的 BPSK 信号双边平均功率谱密度为
\[
P_{BPSK}(f)=\frac{A^2T_b}{4}
\Bigl[
\sinc^2((f-f_c)T_b)+\sinc^2((f+f_c)T_b)
\Bigr].
\]
谱形关于 $0$ 对称，以 $\pm f_c$ 为中心；若以题中比特间隔为 $T_b$，则主瓣第一零点位于
\[
f=\pm\left(f_c\pm \frac{1}{T_b}\right).
\]

\textbf{(2)} 常用载波提取方法有两种：
\begin{enumerate}
\item 平方环法：先对 BPSK 信号平方，去掉 $\pm1$ 造成的相位翻转，再经窄带滤波和分频恢复载波；
\item Costas 环法：利用 I/Q 正交乘法器、低通和环路滤波器构成闭环，锁定本地振荡器相位，实现相干载波恢复。
\end{enumerate}

\medskip
\exercise{习题 6.8}
\textbf{题目：}
假设 BPSK 信号在信道传输中受到 AWGN 干扰，噪声均值为 0，双边功率谱密度
\[
\frac{N_0}{2}=10^{-10}\,\mathrm{W/Hz},
\]
发送信号比特能量
\[
E_b=\frac12A^2T_b,
\]
其中 $T_b$ 为比特间隔，$A$ 为信号幅度。若最佳接收的平均误比特率
\[
P_b=10^{-3},
\]
求输入信息速率 $R_b$ 分别为
\[
10\,\mathrm{kbit/s},\quad 100\,\mathrm{kbit/s},\quad 1\,\mathrm{Mbit/s}
\]
时的 BPSK 信号幅度 $A$。

\par\textbf{解：}
由题给条件
\[
\frac12\operatorname{erfc}\!\left(\sqrt{\frac{E_b}{N_0}}\right)=10^{-3}
\]
可知
\[
\frac{E_b}{N_0}=10^{6.8/10}\approx 4.7863.
\]
又因
\[
N_0=2\times 10^{-10}\,\mathrm{W/Hz},
\qquad
T_b=\frac1{R_b},
\]
故
\[
\frac12A^2T_b=4.7863N_0
\]
即
\[
A=\sqrt{2R_b\cdot 4.7863N_0}.
\]

分别代入三种速率：
\[
R_b=10\,\mathrm{kbit/s}\Rightarrow A\approx 4.376\,\mathrm{mV},
\]
\[
R_b=100\,\mathrm{kbit/s}\Rightarrow A\approx 13.84\,\mathrm{mV},
\]
\[
R_b=1\,\mathrm{Mbit/s}\Rightarrow A\approx 43.76\,\mathrm{mV}.
\]

\medskip
\exercise{习题 6.9}
\textbf{题目：}
一 DPSK 数字通信系统，信息速率为 $R_b$，输入数据为
\[
110100010110\cdots
\]
\begin{enumerate}
\item 写出相对码（设相对码的第一个比特为 1）；
\item 画出 DPSK 发送框图；
\item 写出 DPSK 发送信号的载波相位（设第一个比特的 DPSK 信号载波相位为 0）；
\item 画出 DPSK 信号的平均功率谱密度图（标明频率值，假设“1”“0”等概率出现）。
\end{enumerate}

\par\textbf{解：}
\textbf{(1)} 采用差分编码后，相对码序列为
\[
1011000011011\cdots
\]

\textbf{(2)} 发送框图可表述为：
\[
\text{输入比特流}
\rightarrow
\text{差分编码}
\rightarrow
\text{BPSK 调制}
\rightarrow
\text{DPSK 信号}.
\]

\textbf{(3)} 若相对码比特“1”对应相位 $0$，“0”对应相位 $\pi$，则各码元载波相位依次为
\[
0,\ \pi,\ 0,\ 0,\ \pi,\ \pi,\ \pi,\ \pi,\ 0,\ 0,\ \pi,\ 0,\ 0,\dots
\]

\textbf{(4)} 差分编码不改变“1”“0”等概、独立的统计特性，因此 DPSK 的平均功率谱密度与对应 BPSK 的谱形相同：
\[
P_{DPSK}(f)=\frac{A^2T_b}{4}
\Bigl[
\sinc^2((f-f_c)T_b)+\sinc^2((f+f_c)T_b)
\Bigr].
\]
它关于 $0$ 对称，以 $\pm f_c$ 为中心，第一零点位于
\[
f=\pm\left(f_c\pm \frac{1}{T_b}\right).
\]

\medskip
\exercise{习题 6.10}
\textbf{题目：}
一数字通信系统的收发信机，发送二进制信息速率为 $2\,\mathrm{Mbit/s}$，发送载频为 $800\,\mathrm{MHz}$，要求发射频谱限于
\[
(800\pm 1)\,\mathrm{MHz}.
\]
请设计最佳频带传输系统：
\begin{enumerate}
\item 画出发送原理框图，并画出数字调制器各点的双边功率谱密度图；
\item 写出该最佳频带传输系统的平均误比特率计算公式。
\end{enumerate}

\par\textbf{解：}
总可用带宽为 $2\,\mathrm{MHz}$，要发送 $2\,\mathrm{Mbit/s}$，故所需频带利用率至少为
\[
1\,\mathrm{bit/s/Hz}.
\]
考虑实现复杂度，取四进制调制最合适，因此选用
\[
\text{QPSK}.
\]
若采用滚降系数
\[
\alpha=1,
\]
则 QPSK 的符号速率为
\[
R_s=\frac{R_b}{2}=1\,\mathrm{MBaud},
\]
恰可使通带限制在
\[
800\pm1\,\mathrm{MHz}.
\]

发送结构可表述为：
\[
\text{输入比特流}
\rightarrow
\text{串并变换}
\rightarrow
\text{I/Q 映射为 }\pm1
\rightarrow
\text{脉冲成形}
\rightarrow
\begin{cases}
\times \cos(2\pi f_ct)\\
\times (-\sin(2\pi f_ct))
\end{cases}
\rightarrow
\text{相加}.
\]

其中 A、B 两个基带支路的功率谱密度都以 $0$ 为中心，带宽为 $\pm1\,\mathrm{MHz}$；C 点 QPSK 信号的功率谱密度以
\[
f_c=800\,\mathrm{MHz}
\]
为中心，主要分布在
\[
799\sim 801\,\mathrm{MHz}.
\]

QPSK 的最佳接收平均误比特率与 BPSK 相同：
\[
P_b=\frac12\operatorname{erfc}\!\left(\sqrt{\frac{E_b}{N_0}}\right).
\]

\medskip
\exercise{习题 6.11}
\textbf{题目：}
某三次群（$34.368\,\mathrm{Mbit/s}$）数字微波系统，载波频率为 $6\,\mathrm{GHz}$，信道频带宽度为 $25.776\,\mathrm{MHz}$。请设计在 AWGN 干扰下无码间干扰的 QPSK 调制及解调最佳频带传输系统。

\par\textbf{解：}
QPSK 每个符号携带 2 bit，因此符号速率为
\[
R_s=\frac{34.368}{2}=17.184\,\mathrm{MBaud}.
\]
带通信道宽度与滚降系数满足
\[
W=R_s(1+\alpha).
\]
代入
\[
W=25.776\,\mathrm{MHz}
\]
得
\[
\alpha=\frac{25.776}{17.184}-1=0.5.
\]

因此最佳方案为：
\begin{enumerate}
\item 发射端：串并变换后，I/Q 两路分别映射为 $\pm1$；
\item I、Q 两路都采用滚降系数 $\alpha=0.5$ 的根升余弦（RRC）成形；
\item 分别与 $\cos(2\pi f_ct)$ 和 $-\sin(2\pi f_ct)$ 相乘后相加形成 QPSK；
\item 接收端：正交相干解调，I/Q 两路分别经同样的 RRC 匹配滤波器，再采样判决、并串变换恢复数据。
\end{enumerate}

该系统满足无码间干扰条件，且在 AWGN 下为最佳接收结构。其误比特率公式仍为
\[
P_b=\frac12\operatorname{erfc}\!\left(\sqrt{\frac{E_b}{N_0}}\right).
\]

\medskip
\exercise{习题 6.12}
\textbf{题目：}
已知 OQPSK 调制器的输入二进制序列中两个二进制符号等概率出现且相互独立，信息速率
\[
R_b=2\,\mathrm{Mbit/s}.
\]
\begin{enumerate}
\item 画出 OQPSK 调制器框图（成形滤波器为矩形 NRZ 脉冲）；
\item 若输入数据为 $1110010\cdots$，画出同相支路与正交支路基带信号 $I(t)$、$Q(t)$ 的波形；
\item 画出 OQPSK 调制信号的单边功率谱密度图。
\end{enumerate}

\par\textbf{解：}
\textbf{(1)} OQPSK 调制器结构为：
\[
\text{串并变换}
\rightarrow
\begin{cases}
\text{I 路：映射为 }\pm1 \rightarrow \text{NRZ}\\
\text{Q 路：映射为 }\pm1 \rightarrow \text{NRZ} \rightarrow \text{延迟 }T_b/2
\end{cases}
\rightarrow
\begin{cases}
\times \cos(2\pi f_ct)\\
\times (-\sin(2\pi f_ct))
\end{cases}
\rightarrow
\text{相加}.
\]
由于
\[
R_b=2\,\mathrm{Mbit/s},
\]
故
\[
T_b=0.5\,\mu s,
\qquad
\frac{T_b}{2}=0.25\,\mu s.
\]

\textbf{(2)} 串并变换后
\[
I\text{ 路数据}=1100\cdots,\qquad Q\text{ 路数据}=1010\cdots.
\]
映射为 $\pm1$ 后，
\[
I(t)
=
\begin{cases}
+1,&0\le t<4T_b,\\
-1,&4T_b\le t<8T_b,
\end{cases}
\]
而 $Q(t)$ 在延迟 $T_b/2$ 后按
\[
+1,-1,+1,-1,\dots
\]
交替变化，翻转时刻依次出现在
\[
t=T_b,\ 3T_b,\ 5T_b,\ 7T_b,\dots
\]

\textbf{(3)} OQPSK 的功率谱密度与 QPSK 近似相同，主谱集中在载频 $f_c$ 附近，第一零点约位于
\[
f=f_c\pm 1\,\mathrm{MHz}.
\]
因此其单边功率谱密度图是以 $f_c$ 为中心的主瓣加对称旁瓣。

\medskip
\exercise{习题 6.13}
\textbf{题目：}
4ASK 信号产生框图如题图所示。输入二进制序列经二电平变四电平后得到符号序列
\[
a_n\in\{-3,-1,+1,+3\},
\]
符号间互不相关，四电平等概率出现。发送滤波器冲激响应为矩形 NRZ 脉冲 $g_T(t)$。画出图中 A 点与 B 点的双边功率谱密度图，并标上频率值。

\par\textbf{解：}
A 点信号是四电平 PAM：
\[
x_A(t)=\sum_n a_n g_T(t-nT_s).
\]
由于
\[
\mathbb{E}[a_n]=0,\qquad
\mathbb{E}[a_n^2]=\frac{9+1+1+9}{4}=5,
\]
故 A 点功率谱密度只有连续谱，形状为以 $f=0$ 为中心的 $\sinc^2$ 型低通谱。

B 点为 4ASK 调制信号
\[
x_B(t)=x_A(t)\cos(2\pi f_ct),
\]
其双边功率谱密度为
\[
P_B(f)=\frac14\Bigl[P_A(f-f_c)+P_A(f+f_c)\Bigr].
\]
因此 B 点的双边谱是将 A 点低通谱平移到
\[
f=\pm f_c
\]
附近得到的两个带通信号谱瓣。

\medskip
\exercise{习题 6.14}
\textbf{题目：}
4ASK 系统发送信号等概取自集合
\[
\{-3f_1(t),-f_1(t),f_1(t),3f_1(t)\},
\]
通过 AWGN 信道传输，接收框图如题图所示。求：
\begin{enumerate}
\item 判决器输入噪声方差；
\item 四种发送条件下判决统计量的均值；
\item 最大似然判决域；
\item 平均误符号率。
\end{enumerate}

\par\textbf{解：}
\textbf{(1)} 匹配滤波器或相关器输出噪声方差为
\[
\sigma^2=\frac{N_0}{2}.
\]

\textbf{(2)} 四种发送条件下，判决统计量均值分别为
\[
-3,\quad -1,\quad +1,\quad +3.
\]

\textbf{(3)} 最大似然判决门限取相邻均值中点，因此 4 个判决域依次为
\[
(-\infty,-2),\qquad [-2,0),\qquad [0,2),\qquad [2,\infty).
\]

\textbf{(4)} 由对称性，外侧两点与内侧两点的错误概率分别为
\[
q=\frac12\operatorname{erfc}\!\left(\frac{1}{\sqrt{N_0}}\right),
\qquad
2q.
\]
因此平均误符号率为
\[
P_s=\frac{q+2q+2q+q}{4}
=\frac32 q
=\frac34\operatorname{erfc}\!\left(\frac{1}{\sqrt{N_0}}\right).
\]

\medskip
\exercise{习题 6.15}
\textbf{题目：}
QPSK 信号表示为
\[
s_i(t)=\sqrt{\frac{2}{T_s}}
\cos\!\left(2\pi f_ct+\frac{2\pi(i-1)}{4}\right),
\qquad i=1,2,3,4,\quad 0\le t\le T_s.
\]
两个归一化正交基函数为
\[
f_1(t)=\sqrt{\frac{2}{T_s}}\cos 2\pi f_ct,\qquad
f_2(t)=-\sqrt{\frac{2}{T_s}}\sin 2\pi f_ct.
\]
求：
\begin{enumerate}
\item 星座图及各星座点坐标；
\item 平均符号能量 $E_s$ 与平均比特能量 $E_b$；
\item 若各星座点等概出现，信道单边噪声功率谱密度 $N_0=0.1\,\mathrm{W/Hz}$，求误比特率和误符号率。
\end{enumerate}

\par\textbf{解：}
\textbf{(1)} 四个波形分别为
\[
s_1(t)=f_1(t),\qquad s_2(t)=f_2(t),\qquad s_3(t)=-f_1(t),\qquad s_4(t)=-f_2(t).
\]
因此对应星座点依次为
\[
1,\quad j,\quad -1,\quad -j,
\]
在二维坐标中分别是
\[
(1,0),\ (0,1),\ (-1,0),\ (0,-1).
\]

\textbf{(2)} 每个符号的能量均为 1，因此
\[
E_s=1.
\]
QPSK 每个符号携带 2 bit，故
\[
E_b=\frac{E_s}{2}=\frac12.
\]

\textbf{(3)} 误比特率为
\[
P_b=\frac12\operatorname{erfc}\!\left(\sqrt{\frac{E_b}{N_0}}\right)
=\frac12\operatorname{erfc}(\sqrt{5})
\approx 7.83\times 10^{-4}.
\]
在高信噪比下，QPSK 的误符号率近似为 BER 的 2 倍：
\[
P_s\approx 2P_b\approx 1.57\times 10^{-3}.
\]

\medskip
\exercise{习题 6.16}
\textbf{题目：}
8PSK 及 8QAM 星座图如题图所示，各星座点等概率出现。
\begin{enumerate}
\item 若 8QAM 星座图中星座点之间的最小欧氏距离为 $A$，求使平均符号能量最小的内圆与外圆半径 $a,b$；
\item 若 8PSK 中相邻星座点间距离为 $A$，求圆半径 $r$；
\item 求两类星座图的平均符号能量。
\end{enumerate}

\par\textbf{解：}
\textbf{(1)} 8QAM 的平均符号能量为
\[
E_s=\frac{a^2+b^2}{2}.
\]
内圆相邻点间距为
\[
d_1=\sqrt2\,a,
\]
内外圆最近点间距为
\[
d_2=\sqrt{a^2+b^2-\sqrt2ab}.
\]
要使平均能量最小，必有
\[
d_1=d_2=A.
\]
于是
\[
a=\frac{A}{\sqrt2},
\qquad
b=\frac{1+\sqrt3}{2}A.
\]

\textbf{(2)} 对于 8PSK，
\[
A=2r\sin\frac{\pi}{8},
\]
故
\[
r=\frac{A}{2\sin(\pi/8)}\approx 1.3066A.
\]

\textbf{(3)} 对 8QAM，
\[
E_{s,8QAM}
=\frac12\left(\frac{A^2}{2}+\frac{(1+\sqrt3)^2A^2}{4}\right)
=\frac{3+\sqrt3}{4}A^2
\approx 1.1830A^2.
\]
对 8PSK，
\[
E_{s,8PSK}=r^2\approx 1.707A^2.
\]
因此在相同最小欧氏距离 $A$ 下，8QAM 的平均符号能量比 8PSK 小，约节省
\[
10\log_{10}\frac{1.707}{1.183}\approx 1.59\,\mathrm{dB}.
\]

\medskip
\exercise{习题 6.17}
\textbf{题目：}
题 6.16 中的每个信号矢量都携带 3 个比特的二进制符号。
\begin{enumerate}
\item 若调制器输入信息速率 $R_b=90\,\mathrm{Mbit/s}$，求 8PSK 与 8QAM 的符号速率 $R_s$；
\item 在给定最小欧氏距离 $A$ 的条件下，比较 8PSK 与 8QAM 所需 $E_b/N_0$ 的比值。
\end{enumerate}

\par\textbf{解：}
\textbf{(1)} 每个符号携带 3 bit，因此两者的符号速率相同：
\[
R_s=\frac{R_b}{3}=\frac{90}{3}=30\,\mathrm{MBaud}.
\]

\textbf{(2)} 因为两者每个符号都携带 3 bit，故在给定相同最小距离 $A$ 时，所需
\[
\frac{E_b}{N_0}
\]
之比等于平均符号能量之比：
\[
\frac{(E_b/N_0)_{8PSK}}{(E_b/N_0)_{8QAM}}
=\frac{E_{s,8PSK}}{E_{s,8QAM}}
=\frac{1.707}{1.183}
\approx 1.443.
\]
折合分贝约为
\[
10\log_{10}(1.443)\approx 1.60\,\mathrm{dB}.
\]

\medskip
\exercise{习题 6.18}
\textbf{题目：}
设计一数字通信系统：二进制序列经 MQAM 调制，以 $2400\,\mathrm{Baud}$ 的符号速率在
\[
300\sim3300\,\mathrm{Hz}
\]
的话音频带信道中传输。
\begin{enumerate}
\item 若 $R_b=9600\,\mathrm{bit/s}$，请画出调制框图并说明，并画出图中各点功率谱图；
\item 若 $R_b=14400\,\mathrm{bit/s}$，按上述要求设计。
\end{enumerate}

\par\textbf{解：}
取载频为话音带宽中心频率：
\[
f_c=\frac{300+3300}{2}=1800\,\mathrm{Hz}.
\]
信道带宽为
\[
3000\,\mathrm{Hz}.
\]
又由于符号速率
\[
R_s=2400\,\mathrm{Baud},
\]
升余弦带宽满足
\[
3000=2400(1+\alpha),
\]
故
\[
\alpha=0.25.
\]

\textbf{(1) $R_b=9600\,\mathrm{bit/s}$}\par
每个符号需携带
\[
\frac{9600}{2400}=4\ \text{bit},
\]
故选
\[
16QAM.
\]
调制结构可表述为：
\[
\text{串并变换}
\rightarrow
\text{I/Q 映射为 4 电平}
\rightarrow
\text{RRC 成形 }(\alpha=0.25)
\rightarrow
\begin{cases}
\times \cos(2\pi f_ct)\\
\times (-\sin(2\pi f_ct))
\end{cases}
\rightarrow
\text{相加}.
\]

\textbf{(2) $R_b=14400\,\mathrm{bit/s}$}\par
每个符号需携带
\[
\frac{14400}{2400}=6\ \text{bit},
\]
故选
\[
64QAM.
\]
调制框图与上题相同，只是 I/Q 两路分别映射为 8 电平。

两种情况下，I 路和 Q 路基带功率谱密度都以 0 为中心，支撑范围大致在
\[
-1500\sim1500\,\mathrm{Hz};
\]
QAM 通带功率谱密度则平移到
\[
300\sim3300\,\mathrm{Hz}
\]
内，中心频率为
\[
1800\,\mathrm{Hz}.
\]

\medskip
\exercise{习题 6.19}
\textbf{题目：}
MPSK 信号表示式为
\[
s_i(t)=A\cos\!\left[2\pi f_ct+\frac{2\pi(i-1)}{M}\right],
\qquad i=1,2,\dots,M,\quad 0\le t\le T_s.
\]
请分别写出 QPSK 及 8PSK 信号的解析信号及复包络表达式。

\par\textbf{解：}
一般 MPSK 的解析信号为
\[
z_i(t)=e^{j\left[2\pi f_ct+\frac{2\pi(i-1)}{M}\right]},
\qquad i=1,2,\dots,M.
\]
对应复包络为
\[
s_{i,L}(t)=e^{j\frac{2\pi(i-1)}{M}},
\qquad i=1,2,\dots,M.
\]

\textbf{QPSK}\par
\[
z_i(t)=e^{j\left(2\pi f_ct+\frac{\pi(i-1)}{2}\right)},
\qquad
s_{i,L}(t)=e^{j\frac{\pi(i-1)}{2}},
\qquad i=1,2,3,4.
\]

\textbf{8PSK}\par
\[
z_i(t)=e^{j\left(2\pi f_ct+\frac{\pi(i-1)}{4}\right)},
\qquad
s_{i,L}(t)=e^{j\frac{\pi(i-1)}{4}},
\qquad i=1,2,\dots,8.
\]

\medskip
\exercise{习题 6.20}
\textbf{题目：}
MQAM 信号表示式为
\[
s_i(t)=a_{ic}g_T(t)\cos\omega_ct-a_{is}g_T(t)\sin\omega_ct,
\qquad i=1,2,\dots,M,
\]
其中 $\{a_{ic}\}$ 及 $\{a_{is}\}$ 是一组离散幅度集合，$g_T(t)$ 是基带成形滤波器的冲激响应。请写出矩形星座 16QAM 信号的解析信号及复包络表示式。

\par\textbf{解：}
矩形星座 16QAM 的解析信号为
\[
z_i(t)=\bigl(a_{ic}+ja_{is}\bigr)g_T(t)e^{j\omega_ct},
\qquad i=1,2,\dots,16,
\]
其复包络为
\[
s_L(t)=\bigl(a_{ic}+ja_{is}\bigr)g_T(t).
\]
对矩形 16QAM，
\[
a_{ic},a_{is}\in\{\pm1,\pm3\}.
\]

\medskip
\exercise{习题 6.21}
\textbf{题目：}
为实现 GSM 系统中 $BT_b=0.3$ 的 GMSK 调制器（二进制符号速率 $R_b=270.833\,\mathrm{kbit/s}$）的基带数字化，请选择合适的离散采样速率 $f_s$（每比特抽样点数）及每个样值的量化比特数，并说明如何生成 $\cos\theta(t)$ 与 $\sin\theta(t)$ 查找表。

\par\textbf{解：}
按教材给出的设计结果，可取采样率为
\[
f_s=8R_b,
\]
即每比特取 8 个样点。

量化比特数可取
\[
10\ \text{bit}.
\]

在此基础上，可按所选过采样率离散计算相位函数 $\theta(t)$，再逐点生成
\[
\cos\theta(t),\qquad \sin\theta(t)
\]
的查找表，供数字化 GMSK 调制器调用。

\medskip
